\chapter{Gestione delle sessioni}

Nel capitolo precedente abbiamo realizzato con Node.js la nostra prima
applicazione web, la lista di cose da fare. Le manca ancora una funzionalità
essenziale: più utenti devono poter registrarsi, autenticarsi e gestire
ciascuno la \textbf{propria} lista. Realizzarla richiede un concetto nuovo, il
\term{tracciamento della sessione}.

\section{Il problema}

Torniamo al punto lasciato in sospeso nel primo capitolo: HTTP è
\textbf{stateless}. Una richiesta, di per sé, non contiene alcuna informazione
sulle richieste precedenti.

Come fa il server a sapere che un utente ha già effettuato il login?

\begin{figure}[H]
  \centering
  \begin{tikzpicture}
    % application server con le due rotte, e il database
    \draw[wtserver] (9.0,1.15) rectangle (13.4,-3.05);
    \node[font=\Lato\small\bfseries, text=primary!45!black] at (11.2,0.8)
      {Application Server};
    \node[wtrotta, minimum width=2.8cm, minimum height=1.3cm] at (11.2,-0.25) {/login};
    \node[wtrotta, minimum width=2.8cm, minimum height=1.3cm] at (11.2,-2.05) {/todo};
    \node[wtdatabase] (db) at (15.3,-0.95) {database};
    \draw[{Stealth[length=5pt,width=4pt]}-{Stealth[length=5pt,width=4pt]},
          line width=0.9pt, draw=accent!70!black] (13.4,-0.95) -- (db.west);

    % l'utente e i quattro messaggi
    \pic at (0,-1.07) {utente=naplesyellow!80!orange};
    \node[mcod, below] at (0,-1.1) {Utente};
    \draw[wtreq]  (0.8,0.1)  -- (9.8,0.1);
    \draw[wtrisp] (9.8,-0.6) -- (0.8,-0.6);
    \draw[wtreq]  (0.8,-1.7) -- (9.8,-1.7);
    \draw[wtrisp] (9.8,-2.4) -- (0.8,-2.4);
    \node[wtmsg] at (5.0,0.1)  {1: POST /login usr=luigi\&pwd=pass};
    \node[wtmsg] at (5.0,-0.6) {2: 200 OK};
    \node[wtmsg] at (5.0,-1.7) {3: GET /todo};
    \node[wtmsg] at (5.0,-2.4) {4: \textcolor{merr}{\textbf{?}}};
    \node[mnota, text=merr, anchor=north] at (5.0,-2.5)
      {la richiesta 3 non dice chi la invia: la lista di chi?};
  \end{tikzpicture}
  \caption{Senza stato, la terza richiesta arriva da uno sconosciuto.}
  \label{fig:stateless-problem}
\end{figure}

\dfn{Tracciamento della sessione}{
  Il \term{tracciamento della sessione} (\textit{session tracking}) consiste nel
  mantenere informazioni \textbf{con stato} sulle interazioni di un utente
  attraverso più richieste HTTP. Permette ai server di \guillemotleft riconoscere\guillemotright\ gli
  utenti e le loro azioni, e quindi di adattare le risposte a utenti e
  situazioni diverse.}

\section{Un approccio ingenuo: i cookie}

Un modo che già conosciamo perché il server \guillemotleft memorizzi\guillemotright\ dati fra
richieste diverse è impostare dei \term{cookie}. L'idea è la seguente.

\begin{enumerate}
  \item L'utente invia una richiesta a una pagina dinamica, con nome utente e
        password come parametri.
  \item La pagina verifica se le credenziali sono corrette.
    \begin{itemize}
      \item Se lo sono, il server imposta dei cookie per tenere traccia
            dell'interazione, per esempio
            \codel{Set-Cookie: auth=true; username=luigi}. È una scrittura
            abbreviata: in realtà ogni cookie ha il proprio header
            \code{Set-Cookie}, come si vedrà nel codice.
      \item Se non lo sono, mostra una pagina di errore e non imposta alcun
            cookie.
    \end{itemize}
  \item Le pagine dinamiche che richiedono autenticazione controllano il cookie
        \code{auth}: se è impostato mostrano il contenuto, altrimenti
        reindirizzano alla pagina di login (o rispondono con un errore).
\end{enumerate}

\begin{figure}[H]
  \centering
  \begin{tikzpicture}
    % application server con le due rotte, e il database
    \draw[wtserver] (9.0,1.15) rectangle (13.4,-4.65);
    \node[font=\Lato\small\bfseries, text=primary!45!black] at (11.2,0.8)
      {Application Server};
    \node[wtrotta, minimum width=2.8cm, minimum height=1.3cm] at (11.2,-0.25) {/login};
    \node[wtrotta, minimum width=2.8cm, minimum height=3.0cm] at (11.2,-2.8) {/todo};
    \node[wtdatabase] (db) at (15.3,-1.75) {database};
    \draw[{Stealth[length=5pt,width=4pt]}-{Stealth[length=5pt,width=4pt]},
          line width=0.9pt, draw=accent!70!black] (13.4,-1.75) -- (db.west);

    % utente A: login, poi la lista
    \pic at (0,-0.37) {utente=naplesyellow!80!orange};
    \node[mcod, below] at (0,-0.4) {Utente A};
    \draw[wtreq]  (0.8,0.1)  -- (9.8,0.1);
    \draw[wtrisp] (9.8,-0.6) -- (0.8,-0.6);
    \node[wtmsg] at (5.0,0.1)  {1: POST /login usr=luigi\&pwd=pass};
    \node[wtmsg] at (5.0,-0.6) {2: 200 OK Set-Cookie: auth=true; username=luigi};

    \pic at (0,-2.07) {utente=naplesyellow!80!orange};
    \node[mcod, below] at (0,-2.1) {Utente A};
    \draw[wtreq]  (0.8,-1.6) -- (9.8,-1.6);
    \draw[wtrisp] (9.8,-2.3) -- (0.8,-2.3);
    \node[wtmsg] at (5.0,-1.6) {3: GET /todo Cookie: auth=true; username=luigi};
    \node[wtmsg] at (5.0,-2.3) {4: 200 OK + la lista di luigi};

    % utente B: nessun cookie
    \pic at (0,-3.77) {utente=red-gradient-3!85!black};
    \node[mcod, below] at (0,-3.8) {Utente B};
    \draw[wtreq]  (0.8,-3.3) -- (9.8,-3.3);
    \draw[wtrisp] (9.8,-4.0) -- (0.8,-4.0);
    \node[wtmsg] at (5.0,-3.3) {5: GET /todo (senza cookie)};
    \node[wtmsg] at (5.0,-4.0) {6: \textcolor{merr}{401 Unauthorized}};
  \end{tikzpicture}
  \caption{L'approccio ingenuo: lo stato viaggia nei cookie, e il server si fida
           di quello che il browser gli rimanda.}
  \label{fig:naive-cookies}
\end{figure}

Dopo il login (1), il server risponde impostando i cookie (2); da quel momento
il browser li allega a \textbf{ogni} richiesta allo stesso sito (3), e il server
li usa per capire chi sta chiedendo che cosa (4). Chi non ha i cookie, come
l'utente B, riceve un errore (5, 6). Nello schema la risposta al login è
semplificata in \code{200 OK}: il codice che segue risponde con un
reindirizzamento, ma ciò che conta è l'header \code{Set-Cookie}, che il browser
memorizza qualunque sia il codice di stato.

\subsection{Gestire il login}

Il form di login ha i campi \code{usr} e \code{pwd} e usa il metodo
\code{POST}; il corpo della richiesta si analizza come per il salvataggio degli
elementi della lista. Se le credenziali sono corrette, il server reindirizza
alla homepage e imposta due cookie:

\begin{affianco}[0.5][c]
\begin{lstlisting}[style=js, name=login.js]
if (isAuthenticated) {
  response.writeHead(302, {
    "Location": "/",
    "Set-Cookie": [
      `auth=true; max-age=${60 * 60}`,
      `username=${user.username}; max-age=${60 * 60}`
    ]
  });
  response.end();
}
\end{lstlisting}
\risultato
\begin{lstlisting}[style=http, name=Risposta generata]
HTTP/1.1 302 Found
Location: /
Set-Cookie: auth=true; max-age=3600
Set-Cookie: username=luigi; max-age=3600
\end{lstlisting}
\wtnota{Due header \code{Set-Cookie}: il browser\\ memorizza due cookie distinti.}
\end{affianco}

Uno contiene il nome utente, l'altro è un semplice valore booleano; entrambi
scadono dopo un'ora (\code{max-age} è espresso in secondi).

\info{Perché 302}{
  \code{302} (\textit{Found}) è un \textbf{reindirizzamento}: il browser
  memorizza i cookie ricevuti e poi segue l'header \code{Location} con una
  \code{GET} a \code{/}, portando con sé i cookie appena impostati. È lo schema
  \textit{Post/Redirect/Get} del capitolo precedente: aggiornare la pagina dopo
  il login non reinvia le credenziali. Lì si usava \code{303}
  (\textit{See Other}), il codice più preciso dopo una \code{POST} perché chiede
  esplicitamente di proseguire con una \code{GET}; i browser però trattano
  \code{302} allo stesso modo.}

Se invece le credenziali non sono corrette, si mostra di nuovo la pagina di
login restituendo un codice \code{401}:

\begin{affianco}[0.6][c]
\codicepiccolo
\begin{lstlisting}[style=js, name=login.js]
if (!isAuthenticated) {
  let renderedContent = pug.renderFile("./templates/login.pug",
    { "error": "Autenticazione fallita. Controlla le credenziali." });
  response.writeHead(401, { "Content-Type": "text/html" });
  response.end(renderedContent);
}
\end{lstlisting}
\risultato
\begin{browser}[localhost:3000/login]
  {\Lora\bfseries\LARGE Login}\\[4pt]
  {\color{css-crimson}\small Autenticazione fallita. Controlla le credenziali.}\\[6pt]
  \wtcampo[3.2cm]{luigi}\\[3pt]
  \wtcampo[3.2cm]{\textbullet\textbullet\textbullet\textbullet}\\[4pt]
  \wtbottone{Accedi}
\end{browser}
\wtnota{La pagina di login, di nuovo, con il messaggio\\ d'errore e lo stato 401.}
\end{affianco}

\subsection{Recuperare le informazioni di sessione}

Vogliamo gestire le richieste in modo diverso a seconda che l'utente sia
autenticato o meno: se un utente non autorizzato prova a visitare \code{/todo}
gli serviamo una pagina di errore con codice \code{401}, mentre un utente
autenticato vede la propria lista.

\begin{lstlisting}[style=js, name=auth.js]
function checkUserAuthentication(request) {
  let cookies = parseCookies(request);
  if (cookies.auth) {
    return [true, cookies.username];
  } else {
    return [false, undefined];
  }
}
\end{lstlisting}

Purtroppo, anche per analizzare i cookie siamo da soli. L'header \code{Cookie}
si presenta così: \codel{Cookie: auth=true; username=luigi}.

\begin{affianco}[0.55][c]
\begin{lstlisting}[style=js, name=cookies.js]
function parseCookies(request) {
  const list = {};
  const cookieHeader = request.headers?.cookie;
  if (!cookieHeader) return list;

  cookieHeader.split(`;`).forEach(function (cookie) {
    let [name, ...rest] = cookie.split(`=`);
    name = name?.trim();
    if (!name) return;

    const value = rest.join(`=`).trim();
    if (!value) return;

    list[name] = decodeURIComponent(value);
  });

  return list;
}
\end{lstlisting}
\risultato
\begin{immagine}[Dall'header \code{Cookie} all'oggetto \code{list}.]
\begin{tikzpicture}[
    st/.style={mcod, anchor=west},
    a/.style={-{Stealth[length=4pt,width=3.5pt]}, line width=0.7pt, draw=neutral!45!black},
    lab/.style={mnota, anchor=west, text=accent!70!black}]
  \node[st] at (0,0)     {"auth=true; username=luigi"};
  \node[st] at (0,-0.9)  {["auth=true", " username=luigi"]};
  \node[st] at (0,-1.8)  {name = "username"};
  \node[st] at (0,-2.2)  {value = "luigi"};
  \node[st] at (0,-3.1)  {\{ auth: "true", username: "luigi" \}};
  \draw[a] (0.3,-0.2) -- (0.3,-0.7) node[lab, midway, xshift=4pt]{split(";")};
  \draw[a] (0.3,-1.1) -- (0.3,-1.6) node[lab, midway, xshift=4pt]{split("="), trim};
  \draw[a] (0.3,-2.4) -- (0.3,-2.9) node[lab, midway, xshift=4pt]{decodeURIComponent};
\end{tikzpicture}
\end{immagine}
\end{affianco}

\info{Perché quel \code{rest.join}}{
  Il codice cerca di essere robusto quando il \textbf{valore} di un cookie
  contiene a sua volta un \code{=}. Spezzare su \code{=} produce più di due
  parti, e ricomporre tutto ciò che segue il primo \code{=} è l'unico modo
  corretto. Di norma i cookie sono codificati come URL, ma la specifica non lo
  impone.}

\subsection{Rendere disponibile lo stato ai controller}

Tutte le pagine della nostra applicazione sono influenzate dall'autenticazione:
la barra di navigazione, per esempio, mostrerà \guillemotleft Benvenuto, utente\guillemotright\
invece di \guillemotleft Login\guillemotright. Vogliamo quindi che questa informazione sia
disponibile a tutti i controller.

Un modo è calcolarla una volta sola in \code{handleRequest} e passarla lungo
la catena come oggetto \term{context}.

\begin{lstlisting}[style=js, name=router.js]
function handleRequest(request, response) {
  let [isUserAuthenticated, username] = checkUserAuthentication(request);

  let context = {
    isUserAuthenticated: isUserAuthenticated,
    username: username
  };

  switch (request.url) {
    /* ... altre rotte ... */
    case "/todo":
      handleTodoListRequest(request, response, context); break;
  }
}
\end{lstlisting}

\begin{affianco}[0.6][c]
\codicepiccolo
\begin{lstlisting}[style=js, name=todo-controller.js]
function handleTodoListRequest(request, response, context = {}) {
  if (!context.isUserAuthenticated) {
    handleError(request, response, 401, "Non autorizzato!");
  } else {
    switch (request.method) {
      case "GET":
        handleTodoListRequestGet(request, response, context); break;
      case "POST":
        handleTodoListRequestPost(request, response, context); break;
      default:
        handleError(request, response, 405, "Metodo non supportato");
    }
  }
}
\end{lstlisting}
\risultato
\begin{immagine}[Prima l'autenticazione, poi il metodo.]
\begin{tikzpicture}[
    dec/.style={diamond, aspect=2.2, draw=primary!60!black, fill=primary!8!white,
      line width=0.7pt, font=\Lato\scriptsize, inner sep=1pt},
    esito/.style={rounded corners=3pt, draw=neutral!50!black, fill=white, line width=0.7pt,
      font=\FiraCodeNL\scriptsize, inner sep=3pt},
    a/.style={-{Stealth[length=4pt,width=3.5pt]}, line width=0.7pt, draw=neutral!45!black}]
  \node[dec] (q) at (0,0) {autenticato?};
  \node[esito, draw=merr, text=merr] (e) at (2.6,0) {401};
  \node[dec] (m) at (0,-1.6) {metodo?};
  \node[esito] (g) at (-1.6,-3.0) {\ldots Get()};
  \node[esito] (p) at (0.2,-3.0) {\ldots Post()};
  \node[esito, draw=merr, text=merr] (x) at (1.8,-3.0) {405};
  \draw[a] (q) -- node[mnota, above]{no} (e);
  \draw[a] (q) -- node[mnota, right]{sì} (m);
  \draw[a] (m) -- node[mcod, left, font=\FiraCodeNL\tiny]{GET} (g);
  \draw[a] (m) -- node[mcod, right, font=\FiraCodeNL\tiny]{POST} (p);
  \draw[a] (m) -- node[mnota, right]{altro} (x);
\end{tikzpicture}
\end{immagine}
\end{affianco}

\info{Il codice completo}{
  Il codice completo di questa versione dell'applicazione, con i dati di
  sessione memorizzati nei cookie, è disponibile tra i materiali del corso su
  Teams. Conviene scaricarlo, eseguirlo e seguirne il funzionamento con il
  debugger, come visto nel capitolo 11.}

\section{I problemi dell'approccio a cookie}

L'approccio funziona ed è semplice da implementare. Purtroppo soffre di un
problema \textbf{critico}: i cookie sono interamente controllati dal client, e
manipolarli è banale. Bastano i DevTools: nella scheda \textit{Archiviazione}
di Firefox (\textit{Applicazione} in Chrome) i cookie di un sito sono elencati
in una tabella, e un doppio clic su un valore permette di modificarlo.

\begin{figure}[H]
\centering
\begin{minipage}{15.4cm}
\begin{devtools}{Archiviazione}
  {\Lato\scriptsize\color{neutral!20!black}%
  \noindent\begin{minipage}[t]{3.4cm}\vspace{5pt}
    \hspace*{6pt}{\tiny$\blacktriangleright$}\ Archiviazione cache\par\vspace{3pt}
    \hspace*{6pt}{\tiny$\blacktriangledown$}\ Cookie\par\vspace{3pt}
    \hspace*{15pt}\colorbox{dtblue!12}{http://localhost:3000}\par\vspace{3pt}
    \hspace*{6pt}{\tiny$\blacktriangleright$}\ Indexed DB\par\vspace{3pt}
    \hspace*{6pt}{\tiny$\blacktriangleright$}\ Archiviazione locale\par\vspace{3pt}
    \hspace*{6pt}{\tiny$\blacktriangleright$}\ Archiviazione sessione\par\vspace{5pt}
  \end{minipage}%
  {\color{dtline}\vrule width 0.8pt}%
  \begin{minipage}[t]{\dimexpr\linewidth-3.4cm-0.8pt\relax}\vspace{0pt}
    \colorbox{dtbar}{\makebox[\dimexpr\linewidth-2\fboxsep][l]{%
      \color{neutral!45!black}\quad Filtra elementi\hfill +\quad}}\par
    \noindent\begin{tabular*}{\linewidth}{@{\hspace{8pt}}l@{\extracolsep{\fill}}lllll@{\hspace{8pt}}}
      \rule{0pt}{11pt}\textbf{Nome} & \textbf{Valore} & \textbf{Dominio} &
        \textbf{Percorso} & \textbf{Scadenza / Max-Age} & \textbf{HttpOnly} \\ \hline
      \rule{0pt}{11pt}auth & true & localhost & / & Mon, 27 Oct 2025 18:42:17 GMT & false \\
      username & \fcolorbox{dtblue}{white}{bob\rule[-1.5pt]{0.5pt}{8pt}} &
        localhost & / & Mon, 27 Oct 2025 18:42:17 GMT & false \\
    \end{tabular*}\par\vspace{4pt}
  \end{minipage}}
\end{devtools}
\end{minipage}
\caption{La scheda \textit{Archiviazione} dei DevTools di Firefox: un doppio
         clic sul valore del cookie \code{username}, e \code{luigi} diventa
         \code{bob}. Alla richiesta successiva il server crede di parlare con
         bob.}
\label{fig:cookie-tampering}
\end{figure}

\error{I cookie sono interamente controllati dal client}{
  Un utente può cambiare il valore del proprio cookie \code{username} in
  qualunque altro nome utente, oppure impostarsi da solo un cookie
  \code{auth=true} \textbf{senza mai passare dalla pagina di login}. La nostra
  autenticazione diventa del tutto inutile.

  Il problema della manomissione si può mitigare con i \term{cookie firmati}: il
  server firma i cookie con la propria chiave privata e, finché la chiave resta
  segreta, riconosce quelli manomessi. Resta però il fatto che i dati viaggiano
  sul client.}

I problemi non si fermano alla manomissione. I browser impongono restrizioni
sui cookie per evitare problemi di prestazioni:

\begin{itemize}
  \item la dimensione di un cookie è generalmente limitata a \textbf{4096 byte};
  \item ogni host può memorizzare in genere dai \textbf{20 ai 50 cookie} per
        dominio.
\end{itemize}

I limiti del proprio browser si possono verificare su
\urlx{http://browsercookielimits.iain.guru/}. E se dovessimo memorizzare più
dati di sessione?

\section{Le sessioni web}

Le domande sono quindi due. Come tenere traccia di informazioni
\textbf{sensibili} sulle interazioni precedenti, senza che il client possa
manometterle? E come memorizzare \textbf{più dati} di quanti ne permettano i
cookie? La risposta è il meccanismo delle \term{sessioni web}: i dati restano
sul server, e il cookie (lo stesso meccanismo della RFC 6265, \textit{HTTP
State Management Mechanism}) trasporta soltanto un identificatore.

\dfn{Sessione web}{
  Quando serve tenere traccia di informazioni attraverso più richieste, il server
  crea una \term{sessione} per un dato client. Una sessione è una struttura dati
  che memorizza coppie chiave-valore, \textbf{gestita e conservata sul server}.
  È identificata da un \term{identificatore di sessione} univoco (o
  \textit{session token}), che viene passato al client, generalmente come
  cookie.}

Il punto chiave è che i \textbf{dati restano sul server}: il client possiede
soltanto un identificatore opaco, e non può metterci mano.

\begin{figure}[H]
  \centering
  \colorlet{sidA}{orange-gradient-6!80!black}%
  \begin{tikzpicture}
    % application server con i dati di sessione
    \draw[wtserver] (9.0,1.15) rectangle (15.2,-4.65);
    \node[font=\Lato\small\bfseries, text=primary!45!black] at (12.1,0.8)
      {Application Server};
    \draw[rounded corners=4pt, draw=accent!55!black, line width=0.8pt,
          fill=accent!6!white] (9.35,0.4) rectangle (14.85,-4.3);
    \node[font=\Lato\scriptsize\bfseries, text=accent!45!black, anchor=west]
      at (9.5,0.12) {Dati di sessione};
    \oggetto[fill=orange-gradient-1!70!white, text=sidA,
             font=\FiraCodeNL\scriptsize\bfseries]
      {s1}{(12.1,-0.5)}{5.0cm}{id: ahs37gfw\ldots asi3ff435}
      {auth/true, username/"john"}
    \oggetto[font=\FiraCodeNL\scriptsize\bfseries]
      {s2}{(12.1,-2.55)}{5.0cm}{id: 4fg23g54\ldots ssdf23g43u}
      {auth/true, username/"billy"}

    % utente A: login, poi la lista
    \pic at (0,-0.37) {utente=naplesyellow!80!orange};
    \node[mcod, below] at (0,-0.4) {Utente A};
    \draw[wtreq]  (0.8,0.1)  -- (9.0,0.1);
    \draw[wtrisp] (9.0,-0.6) -- (0.8,-0.6);
    \node[wtmsg] at (4.9,0.1)  {1: POST /login usr=john\&pwd=pass};
    \node[wtmsg] at (4.9,-0.6)
      {2: 200 OK Set-Cookie: SID=\textcolor{sidA}{ahs37\ldots 435}};

    \pic at (0,-2.07) {utente=naplesyellow!80!orange};
    \node[mcod, below] at (0,-2.1) {Utente A};
    \draw[wtreq]  (0.8,-1.6) -- (9.0,-1.6);
    \draw[wtrisp] (9.0,-2.3) -- (0.8,-2.3);
    \node[wtmsg] at (4.9,-1.6)
      {3: GET /todo Cookie: SID=\textcolor{sidA}{ahs37\ldots 435}};
    \node[wtmsg] at (4.9,-2.3) {4: 200 OK + la lista di john};

    % utente B: nessun identificatore
    \pic at (0,-3.77) {utente=red-gradient-3!85!black};
    \node[mcod, below] at (0,-3.8) {Utente B};
    \draw[wtreq]  (0.8,-3.3) -- (9.0,-3.3);
    \draw[wtrisp] (9.0,-4.0) -- (0.8,-4.0);
    \node[wtmsg] at (4.9,-3.3) {5: GET /todo (senza SID)};
    \node[wtmsg] at (4.9,-4.0) {6: redirect alla pagina di login};
  \end{tikzpicture}
  \caption{Il client possiede solo l'identificatore; i dati restano sul server.
           Chi non presenta un identificatore valido viene rimandato al login.}
  \label{fig:sessions}
\end{figure}

Seguiamo i messaggi dello schema.
\begin{itemize}[leftmargin=2.6em, labelwidth=2em, labelsep=0.6em, align=left]
  \item[\textbf{1-2}] L'utente A invia le credenziali di john. Il server crea
        una nuova sessione, vi memorizza \code{auth} e \code{username} e
        restituisce al browser \textbf{solo} l'identificatore, nel cookie
        \code{SID} (nel codice che segue si chiamerà \code{sessionId}).
  \item[\textbf{3-4}] A ogni richiesta successiva il browser rimanda il cookie.
        Il server cerca la sessione con quell'identificatore, vi trova
        \code{username: "john"} e mostra la lista di john.
  \item[\textbf{5-6}] L'utente B non ha alcun identificatore valido: la sua
        richiesta non corrisponde a nessuna sessione, e il server lo
        reindirizza alla pagina di login.
\end{itemize}

\section{Implementare le sessioni da zero}

Realizzare un meccanismo di sessione di base in Node.js è piuttosto semplice.
Si possono usare gli oggetti \code{Map} di JavaScript: una prima mappa associa
ciascun identificatore di sessione a una sessione, che è a sua volta una mappa
di coppie chiave-valore.

\begin{affianco}[0.56][c]
\begin{lstlisting}[style=js, name=Session.js]
class Session {

  constructor() {
    this.sessionStore = new Map();
  }

  createSession() {
    let sessionId = this.generateNewSessionId();
    this.sessionStore.set(sessionId, new Map());
    return sessionId;
  }

  storeSessionData(sessionId, key, value) {
    return this.getSessionById(sessionId)?.set(key, value);
  }

  getSessionFromRequest(request) {
    let requestSessionId = parseCookies(request)["sessionId"];
    return this.getSessionById(requestSessionId);
  }

  /* altri metodi ... */
}
\end{lstlisting}
\risultato
\begin{immagine}[Una mappa di mappe: a ogni identificatore\\
                 corrisponde una sessione.]
\begin{tikzpicture}
  \mchiave{1.6cm}
  \oggetto[fill=accent!14!white]{st}{(0,-0.35)}{2.8cm}{sessionStore (Map)}
    {Kjhdfgo7\ldots 34niAdudh/, 34uih\&7G\ldots Q4jkhsadg/,
     908shdu\%\ldots jhsdf\$a\&a/}
  \mchiave{1.4cm}
  \oggetto[fill=orange-gradient-1!70!white]{s1}{(4.75,1.0)}{2.6cm}{Map}
    {username/"luigi", auth/true}
  \oggetto{s2}{(4.75,-0.6)}{2.6cm}{Map}{username/"bob", auth/true}
  \oggetto[fill=green-gradient-1]{s3}{(4.75,-2.2)}{2.6cm}{Map}
    {username/"bob", auth/true}
  \node (d1) at ([xshift=-9pt]st-1.east) {};
  \node (d2) at ([xshift=-9pt]st-2.east) {};
  \node (d3) at ([xshift=-9pt]st-3.east) {};
  \riferimento{d1}{to[out=0, in=180] (s1.west)}
  \riferimento{d2}{to[out=0, in=180] (s2.west)}
  \riferimento{d3}{to[out=0, in=180] (s3.west)}
\end{tikzpicture}
\end{immagine}
\end{affianco}

Le chiavi della prima mappa sono gli identificatori di sessione, i valori le
sessioni. Nell'esempio bob ha \textbf{due} sessioni: è normale, perché
ogni login ne crea una nuova. Basta che bob abbia effettuato l'accesso da due
browser, o da due dispositivi diversi, e ciascuno avrà il proprio
identificatore.

Il listato omette i metodi rimanenti: la versione completa è nel file
\code{Session.js} dell'esempio del corso. Una possibile implementazione dei due
metodi usati sopra è la seguente.

\begin{affianco}[0.46][c]
\begin{lstlisting}[style=js, name=Session.js]
import { randomUUID } from "crypto";

class Session {
  /* ... i metodi visti sopra ... */

  generateNewSessionId() {
    // generatore crittograficamente sicuro
    return randomUUID();
  }

  getSessionById(sessionId) {
    // undefined se la sessione non esiste
    return this.sessionStore.get(sessionId);
  }
}
\end{lstlisting}
\risultato
\begin{immagine}[Un identificatore prevedibile si indovina; uno casuale no.]
\begin{tikzpicture}[
    a/.style={-{Stealth[length=4pt,width=3.5pt]}, line width=0.7pt},
    sid/.style={rounded corners=2pt, fill=white, line width=0.7pt,
      font=\FiraCodeNL\scriptsize, inner sep=3pt}]
  \node[mnota, anchor=west] at (0,0.5) {\textbf{con un contatore}};
  \node[sid, draw=neutral!50!black, anchor=west] (a) at (0,0) {sessionId=41};
  \node[sid, draw=merr, text=merr, anchor=west] (b) at (3.2,0) {sessionId=42};
  \draw[a, draw=merr] (a) -- node[mnota, above, text=merr]{+1} (b);
  \node[mnota, text=merr, anchor=north west] at (0,-0.3)
    {basta provare il valore successivo per rubare la sessione di un altro};
  \node[mnota, anchor=west] at (0,-1.25) {\textbf{con randomUUID()}};
  \node[sid, draw=mok, anchor=west] at (0,-1.75)
    {sessionId=1b9d6bcd-bbfd-4b2d-9b5d-ab8dfbbd4bed};
  \node[mnota, text=mok, anchor=north west] at (0,-2.05)
    {122 bit casuali: indovinarne uno valido è praticamente impossibile};
\end{tikzpicture}
\end{immagine}
\end{affianco}

\warning{L'identificatore deve essere imprevedibile}{
  Chi conosce un identificatore valido \textbf{è}, per il server, il
  proprietario della sessione. Se gli identificatori fossero numeri progressivi,
  o generati con \code{Math.random()}, che non è pensato per la sicurezza, un
  attaccante potrebbe indovinarli. Per questo si usa un generatore
  crittograficamente sicuro, come \code{randomUUID} del modulo \code{crypto}.
  Il furto della sessione (\textit{session hijacking}) e le sue contromisure
  sono trattati nel capitolo 23.}

\subsection{Usare le sessioni}

Quando un utente effettua il login, creiamo una nuova sessione, vi memorizziamo
i dati rilevanti e passiamo l'identificatore come cookie.

\begin{affianco}[0.5][c]
\codicepiccolo
\begin{lstlisting}[style=js, name=login.js]
if (isAuthenticated) {
  // crea una nuova sessione per l'utente
  let sessionId = session.createSession();
  session.storeSessionData(sessionId, "username", user.username);
  session.storeSessionData(sessionId, "auth", true);

  response.writeHead(302, {
    "Location": "/",
    "Set-Cookie": [
      `sessionId=${sessionId}; max-age=${60 * 60}`
    ]
  });
  response.end();
}
\end{lstlisting}
\risultato
\begin{lstlisting}[style=http, name=Risposta generata]
HTTP/1.1 302 Found
Location: /
Set-Cookie: sessionId=ahs37...; max-age=3600
\end{lstlisting}
\wtnota{Al client arriva solo l'identificatore:\\
        nome utente e \code{auth} restano sul server.}
\end{affianco}

Quando serve accedere ai dati di sessione:

\begin{affianco}[0.55][c]
\begin{lstlisting}[style=js, name=router.js]
function handleRequest(request, response) {
  let userSession = session.getSessionFromRequest(request);

  let context = {
    isUserAuthenticated: userSession?.get("auth"),
    username: userSession?.get("username")
  };

  /* resto del codice */
}
\end{lstlisting}
\risultato
\begin{immagine}[Dal cookie alla sessione, dalla sessione al contesto.]
\begin{tikzpicture}[
    st/.style={rounded corners=3pt, draw=primary!55!black, fill=primary!6!white,
      line width=0.7pt, font=\FiraCodeNL\scriptsize, inner sep=3pt, align=center},
    a/.style={-{Stealth[length=4pt,width=3.5pt]}, line width=0.7pt, draw=neutral!45!black}]
  \node[st] (c) at (0,0) {Cookie: sessionId=ahs37\ldots};
  \node[st, fill=accent!12!white, draw=accent!55!black] (s) at (0,-1.1) {sessionStore.get("ahs37\ldots")};
  \node[st] (u) at (0,-2.2) {userSession\\ username: "john", auth: true};
  \node[st, fill=green-gradient-1!45!white, draw=green-gradient-6] (x) at (0,-3.4)
    {context\\ isUserAuthenticated: true};
  \draw[a] (c) -- (s); \draw[a] (s) -- (u); \draw[a] (u) -- (x);
\end{tikzpicture}
\end{immagine}
\end{affianco}

Il codice del router è praticamente identico a quello della versione con i
cookie: cambia solo \textbf{da dove} arrivano i dati. Il cookie ora contiene un
identificatore opaco, e manometterlo non serve a nulla: se l'identificatore non
corrisponde a una sessione esistente sul server, \code{getSessionById}
restituisce \code{undefined}.

\warning{Le sessioni in memoria non bastano in produzione}{
  Conservare le sessioni in una \code{Map} in memoria funziona benissimo per
  imparare, ma ha due limiti concreti. Al \textbf{riavvio} del server tutte le
  sessioni spariscono e tutti gli utenti devono rifare il login. E se
  l'applicazione gira su \textbf{più macchine} dietro un bilanciatore di carico,
  una richiesta servita dalla macchina B non trova la sessione creata sulla
  macchina A. In produzione le sessioni si conservano in un archivio condiviso
  (Redis, un database), oppure si adottano token autoconsistenti e firmati.}

\info{Il codice completo, versione con le sessioni}{
  Anche la versione migliorata dell'applicazione, con le sessioni memorizzate
  sul server e la classe \code{Session} completa, è disponibile tra i materiali
  del corso su Teams. Vale lo stesso consiglio: eseguirla e seguirla con il
  debugger, confrontandola con la versione a cookie. Nel capitolo 15 lo stesso
  meccanismo sarà offerto già pronto dal middleware \code{express-session}.}

\subsection{Altri metodi di tracciamento}

Usare un cookie per memorizzare l'identificatore di sessione è l'approccio più
diffuso: è semplice e ampiamente supportato. Ne esistono però altri:

\begin{itemize}
  \item \term{URL rewriting}: l'identificatore viene aggiunto a ogni URL come
        parametro di query;
  \item \term{campi form nascosti}: l'identificatore è memorizzato in campi
        nascosti dei form e viene inviato insieme agli altri dati.
\end{itemize}

Entrambi hanno lo svantaggio di esporre l'identificatore in luoghi dove
è facile che finisca per essere condiviso o registrato nei log: si usano
soprattutto quando i cookie non sono disponibili.

\section{Riferimenti}

\begin{itemize}
  \item \textbf{Introduction to Web Applications Development}, Carles Mateu.
        Liberamente disponibile su \textit{archive.org} sotto licenza GNU FDL. \\
        Parte rilevante: sezione 3.11 (\textit{Session monitoring}); si possono
        ignorare gli esempi con le Java Servlet.
  \item \textbf{Using HTTP cookies} --- MDN web docs. \\
        \urlx{https://developer.mozilla.org/en-US/docs/Web/HTTP/Cookies}
  \item \textbf{RFC 6265: HTTP State Management Mechanism}. \\
        \urlx{https://datatracker.ietf.org/doc/html/rfc6265}
\end{itemize}
